\section{Introduction}
\label{sec:introduction}

The discovery and design of novel materials is increasingly framed as a sequential decision-making problem: compute or synthesise a candidate, measure its properties, learn from the outcome, and suggest the next candidate~\cite{frazier2018tutorial}.
Bayesian optimisation (BO) has emerged as the engine of this design loop, offering a principled framework for balancing the exploration of unknown chemical space against the exploitation of promising regions~\cite{jones1998efficient,todorovic2019bayesian}.
At the heart of any BO campaign lies a \emph{surrogate model}---typically a Gaussian process (GP)---that approximates the expensive property-evaluation function and quantifies predictive uncertainty to guide acquisition.
However, the effectiveness of a surrogate depends critically on two upstream choices: \emph{what features represent the crystal structure}, and \emph{what probabilistic model maps those features to properties}.
These two choices interact combinatorially, yet no systematic study has evaluated their joint effect across the diversity of material property domains that real design campaigns must navigate.

That diversity is substantial.
Materials design targets span mechanical properties such as elastic moduli and Poisson's ratio, where the structure--property relationship is dominated by local bonding geometry; electronic properties such as band gap and dielectric constants, where orbital hybridisation and many-body screening play essential roles; and vibrational properties such as phonon spectra and phonon-derived dielectric response, which depend on the potential energy surface beyond the harmonic regime.
A surrogate pipeline that excels on bulk modulus prediction may fail entirely on dielectric constants---not because the model is poorly trained, but because the feature representation lacks the information content needed for the target physical phenomenon.
Understanding which featuriser--surrogate combinations transfer across this property spectrum is therefore a prerequisite for deploying BO in realistic, multi-objective materials design campaigns.

On the featurisation side, the materials science community has historically relied on hand-crafted structural descriptors.
The Smooth Overlap of Atomic Positions (SOAP)~\cite{bartok2013representing} encodes local atomic environments as power spectra of smooth densities, atom-centred symmetry functions (ACSFs)~\cite{behler2011atom} provide another local perspective, and compositional feature sets like Magpie~\cite{ward2016general} and matminer~\cite{ward2018matminer} aggregate element-level statistics.
While effective within their design scope, these descriptors share fundamental limitations: they are hand-crafted with fixed functional forms, operate within fixed receptive fields (typically 4--8~\AA{} cutoffs), and encode no transferable pre-training knowledge.
Each new materials system or target property may require descriptor redesign, hyperparameter tuning, or domain-specific feature engineering---a bottleneck that compounds across the diverse property domains described above.

The emergence of foundation models for atomistic simulation represents a paradigm shift.
Models such as MACE~\cite{batatia2022mace,batatia2024foundation}, ORB~\cite{orb2024}, and UMA/eSEN~\cite{uma2025} have been pre-trained on millions of density functional theory (DFT) calculations spanning broad swaths of chemical space.
Their learned representations encode rich structural and chemical information---not just local bond geometry, but also long-range ordering, charge transfer tendencies, and composition-dependent energy landscapes.
The key insight is that these representations can be \emph{repurposed}: just as ImageNet-trained convolutional neural networks produce features that transfer to medical imaging or satellite classification, materials foundation models produce embeddings that may serve as general-purpose featurisers for downstream property prediction.
The appeal is immediate: instead of designing a new descriptor for each property domain, one simply extracts embeddings from a pre-trained model and feeds them to a surrogate.
Whether this promise holds across the full spectrum of material properties---from mechanical to electronic to vibrational---is an open question.

On the surrogate side, the Bayesian optimisation community has developed several probabilistic modelling approaches beyond the standard single-task GP~\cite{rasmussen2006gaussian}.
Multi-task GPs (MTGPs)~\cite{bonilla2007multi,swersky2013multi} extend the framework to jointly model multiple correlated properties, sharing statistical strength across tasks through the intrinsic coregionalisation model (ICM) kernel.
Deep GPs (DGPs)~\cite{damianou2013deep,dutordoir2021deep} stack multiple GP layers to capture non-linear input--output relationships that a single-layer GP cannot represent---a capability that may prove essential for properties with complex, non-linear dependence on structure.
The BoTorch framework~\cite{balandat2020botorch}, built on GPyTorch~\cite{gardner2018gpytorch}, provides unified implementations of all three surrogate families, enabling fair comparison under identical training and evaluation protocols.
Each surrogate brings distinct trade-offs: the GP offers stability and interpretability, the DGP offers expressive power at the cost of configuration sensitivity, and the MTGP offers inter-task transfer that may compensate for weak individual predictions.

A further subtlety, often overlooked in surrogate benchmarking, concerns \emph{what metric we use to judge surrogate quality}.
The coefficient of determination (\Rtwo{}) is the standard regression metric, but BO acquisition functions---Expected Improvement, Upper Confidence Bound, Thompson Sampling---depend fundamentally on the surrogate's ability to \emph{rank} candidates relative to the current best, not to predict their exact property values.
A surrogate that correctly orders material A above material B will guide the acquisition function to evaluate A first, regardless of whether it predicts A's property as 150~GPa or 180~GPa.
This observation suggests that Spearman rank correlation ($\rho$) may be a more informative metric for evaluating BO surrogates than \Rtwo{}---a hypothesis that, as we will show, has profound consequences for how we assess surrogate utility across different property domains.

Despite the rapid progress on both fronts---foundation model embeddings and advanced surrogates---\emph{no systematic study has evaluated the cross-product of these choices across diverse material property domains}.
Existing works typically evaluate one featuriser with one surrogate~\cite{batatia2024foundation}, or benchmark multiple models but only with fixed descriptors~\cite{dunn2020benchmarking}.
The practical materials scientist faces a combinatorial design space of featuriser $\times$ surrogate $\times$ dimensionality $\times$ training budget $\times$ property type, with no clear guidance on which combinations work best under which conditions.

In this paper, we fill this gap with the largest cross-factorial benchmark to date for probabilistic surrogate evaluation in materials science.
Our contributions are fourfold:
\begin{enumerate}[nosep]
    \item \textbf{First systematic combinatorial benchmark across diverse property domains.}
    We evaluate 4~featurisers (SOAP, MACE, ORB, UMA/eSEN) $\times$ 3~surrogates (GP, MTGP, DGP) $\times$ 3~PCA dimensionalities $\times$ 3~training sizes across 3~materials property datasets---elastic tensor (8~mechanical properties), dielectric constant (4~electronic properties), and phonon dielectric (3~vibrational/electronic properties)---totalling over 8{,}640 cross-validated experiments spanning 15~target properties.
    \item \textbf{Ranking accuracy as a primary evaluation metric.}
    We discover that Spearman rank correlation ($\rho$) systematically and substantially exceeds \Rtwo{} across all datasets, with the gap widening as properties become more difficult to predict.
    Surrogates with $\bar{R}^2 = 0.303$ can still achieve $\bar{\rho} = 0.824$, and the most extreme case---eps\_electronic with $R^2 = -0.001$ but $\rho = 0.930$---demonstrates that surrogates dismissed by conventional regression metrics can be highly effective for BO.
    We argue that Spearman correlation should be adopted as a primary metric for BO surrogate evaluation.
    \item \textbf{ORB embeddings as the universal default featuriser.}
    ORB outperforms all alternatives on every dataset and every surrogate, achieving $\bar{R}^2 / \bar{\rho}$ of $0.803 / 0.842$ (elastic), $0.336 / 0.778$ (dielectric), and $0.303 / 0.824$ (phonon) at its best configurations.
    This universality stems from ORB's large pre-training corpus ($\sim$120M configurations) and its graph attention architecture optimised for periodic crystals.
    \item \textbf{Practical pipeline configuration guidelines.}
    We characterise the DGP as a powerful surrogate when properly configured---achieving the best \Rtwo{} on both elastic ($\bar{R}^2 = 0.807$) and dielectric ($\bar{R}^2 = 0.336$) datasets---but one that requires PCA dimensionality as a carefully tuned hyperparameter (PCA\,$\leq$\,25).
    The GP remains the robust default, and the MTGP helps selectively when tasks are few, correlated, and individually hard to predict.
\end{enumerate}

We release the complete pipeline as open-source software, providing a community benchmark for future featuriser and surrogate development in materials design.
